% VI/42: Supplementary reader bridge. Original problems and proofs remain intact.
\paragraph{From local generators to a line bundle.}
An invertible $\mathcal O_X$-module $\mathcal L$ is a sheaf that is
locally free of rank one: every point admits a neighborhood $U_i$ on
which $\mathcal L|_{U_i}\simeq\mathcal O_{U_i}$.  Choose a local
basis $e_i$ on each $U_i$.  On an overlap there is a unique unit
$g_{ij}\in\mathcal O_X^\times(U_i\cap U_j)$ such that
$e_i=g_{ij}e_j$.  The identities
\[
 g_{ii}=1,\qquad g_{ij}g_{jk}=g_{ik}
\]
express precisely the consistency needed to glue the local free
modules into one sheaf.  A line bundle can be locally trivial
without admitting a \emph{global} nowhere-vanishing section.
A global trivialization exists exactly when a global generating
section exists.

\paragraph{Worked example 1: the transition for $\mathcal O_{\mathbb P^1}(1)$.}
Let $\mathbb P^1_k=\operatorname{Proj}k[X,Y]$ and cover it by
$U_Y=D_+(Y)$ and $U_X=D_+(X)$.  On $U_Y$ use $t=X/Y$;
on $U_X$ use $s=Y/X=1/t$.  The line bundle
$\mathcal O_{\mathbb P^1}(1)$ has local frames $e_Y=Y$ and
$e_X=X$, where these symbols denote degree-one homogeneous
expressions in the usual chartwise construction.  On the overlap,
\[
 e_X=t\,e_Y,\qquad e_Y=s\,e_X.
\]
The transition is a unit because $t$ is invertible exactly where
both coordinates are nonzero.  Thus the fibers are one-dimensional,
but the local bases are related by a nonconstant function.

A homogeneous linear form $aX+bY$ gives a global section: in the
$Y$-frame it is $(at+b)e_Y$, and in the $X$-frame it is
$(a+bs)e_X$.  Indeed $(a+bs)e_X=(at+b)e_Y$ on the overlap.
Consequently
\[
 H^0(\mathbb P^1_k,\mathcal O(1))
  =kX\oplus kY,\qquad \dim_k H^0(\mathbb P^1_k,\mathcal O(1))=2.
\]
This is an explicit gluing calculation, not merely a count of
homogeneous polynomials.

\paragraph{Worked example 2: twists, duals, and section counts.}
Tensoring the transition $e_X=t e_Y$ with itself $n$ times
produces $\mathcal O(n)$ with transition $e_X^{(n)}=t^n e_Y^{(n)}$
for every integer $n$.  For $n<0$ the notation means the
appropriate tensor power of the dual line bundle.  A section of
$\mathcal O(n)$ with $n\geq0$ is represented on $U_Y$ by a
polynomial $a_0+a_1t+\cdots+a_nt^n$, whose expression on $U_X$
is $s^n(a_0+a_1s^{-1}+\cdots+a_ns^{-n})$.  This is regular
at $s=0$ exactly because its degree in $t$ does not exceed $n$.
Thus
\[
 h^0\!\left(\mathbb P^1,\mathcal O(n)\right)=
 \begin{cases}n+1,&n\geq0,\\0,&n<0.\end{cases}
\]
For negative $n$, the same transition calculation forces a
would-be section to vanish; no nonzero Laurent polynomial can
be regular on both charts with the required twist.
In particular, $\mathcal O(-1)$ is not isomorphic to
$\mathcal O(1)$, despite both being locally free of rank one.

\paragraph{Worked example 3: why $\operatorname{Pic}(\mathbb P^1_k)=\mathbb Z$.}
The Picard group consists of isomorphism classes of line bundles,
with tensor product as its group law and $\mathcal O_X$ as its
identity.  On the smooth curve $\mathbb P^1_k$, Cartier divisors
and line bundles correspond by $D\mapsto\mathcal O(D)$.
The computation in VI/41 shows that every divisor class is
uniquely determined by its degree, using the rational point at
infinity as generator.  Therefore
\[
 \operatorname{Pic}(\mathbb P^1_k)\simeq\mathbb Z,
 \qquad [\mathcal O(n)]\longleftrightarrow n.
\]
The group operation really is addition of integers because
$\mathcal O(a)\otimes\mathcal O(b)\simeq\mathcal O(a+b)$:
on the overlap their transition functions multiply to $t^{a+b}$.
In contrast, $\operatorname{Pic}(\mathbb A^1_k)=0$ because
$k[t]$ is a principal ideal domain and every invertible module
is free.  The distinction is geometric, not a difference in rank.

\paragraph{Connection to Cartier and Weil divisor classes.}
On a locally factorial integral scheme, every Weil divisor is
Cartier, so the Picard group and the Weil class group agree.
On singular normal schemes, however, some Weil divisors fail to
be Cartier: the nontrivial class on the quadric cone from VI/40--VI/41
need not arise from a line bundle.  Thus an equality of Picard
and Weil class groups must never be asserted without hypotheses.
